\section{Conclusion and future work}
\label{sec:conclusion}

Small tool-using agents with persistent memory fail in a particular way when their records disagree: they follow a
reliability label instead of the text, and enforcing that label at the terminal action makes the failure worse.
This paper tested a specific alternative --- arbitration that states the \emph{evidence basis} for each resolution
and enforces only what that basis supports --- under a pre-registration frozen before the confirmatory data
existed.

The pre-registered result is positive and stable. Across 57{,}888 runs on 268 sealed tasks, gated evidence-basis
arbitration beat enforced record-and-temporal labels by $17.1$\pp{} of task success, advisory temporal labels by
$14.5$\pp, an LLM arbiter under the same policy by $15.9$\pp, and unlabelled memory by $15.3$\pp, all Holm
corrected, in the same direction for all six models and stable across three item-set definitions. Harmful answers
fell from 22.6\% to 3.1\%, and false resolution on genuinely unresolvable groups from 19.0\% to 1.9\%.

Three findings qualify what that means, and we regard them as the more useful part of the paper. Most of the gain
is in the labels, not the enforcement: evidence-basis labels alone are worth $11.8$\pp{} over metadata labels,
with gating adding $6.0$\pp. Ungated enforcement of those same labels scores $9.7$\pp{} higher still while doubling
harm and restoring the arbitrary picks the mechanism was built to prevent --- so gating trades accuracy for
safety, and this benchmark does not price that trade. And the same gating driven by an LLM arbiter produced no
gain and more harm, which suggests the general lesson is not about gating but about what enforcement is allowed to
enforce.

The best condition still answers only 61.3\% of tasks correctly, and its residual failure is over-conservatism:
20.2\% false conflicts and 11.5\% unnecessary abstentions against 3.1\% harm. The mechanism moved errors where it
was designed to move them. Whether that is the right place for them is a deployment judgement this study does not
make.

\paragraph{Future work.} Four directions follow directly. (1) A benchmark that assigns explicit costs to a wrong
confident answer and to a missed correct one, so that the gating trade can be decided rather than described.
(2) Extending the arbiter beyond lexical English cues and ISO value dates, where its documented limitations
currently live, and testing it on a set that excludes explicitly signalled invalidation
\citep{chao2026stale}. (3) A transcript study of the remaining false conflicts, to separate items the arbiter
could have resolved from items where abstaining is correct; this needs the data release. (4) Model-family
coverage: the Llama results show the mechanism helping by a different route, and a family whose dominant failure
is abstention needs a different account of what reliability means.

\paragraph{Two methodological artefacts} come out of the campaign rather than the results, and we release both.
Run manifests should record the host kernel, GPU driver and boot identifier, not only framework versions: in this
campaign a mid-campaign platform change was visible only as a 0.2\,GiB difference in a telemetry field. And an
integrity gate should separate what \emph{identifies} a run from what merely \emph{observes} its host, fail closed
on anything unclassified, and report the differences it is willing to tolerate rather than silently accepting or
silently rejecting them.
